\part{The idea}

\chapter{What \nibb{} is}

\nibb{} is a coding agent's \emph{harness}: the program that sits between a
language model and your machine. It holds the conversation, offers the model
tools, runs the tools the model asks for, and feeds the results back. Every
such harness answers three questions, and \nibb{} exists because it answers
them differently.

\paragraph{Who decides what the model may do?} In \nibb{}, one function. Every
tool the model can call is a \emph{primitive} in an embedded Scheme
interpreter, and every primitive application, whoever asks for it, passes
through a single Go function where a \emph{guard} decides: allow, ask the
user, or deny. There is no second path. The Scheme interpreter is the only
surface through which anything reaches the host, and the guard sits on it.

\paragraph{How are secrets protected?} By never being content. A key is not
a string the harness tries to spot in a request; it is a \emph{credential}
the user registered into an encrypted store, referenced by an opaque handle,
and injected only at a destination the user has \emph{bound} it to. The
question ``is this egress safe?'' is not ``does this look like a secret?''
but ``is this class of credential authorised for this destination, in this
position, with this method?''\,---\,a table lookup, owned by the user. Content
detection exists too (a cheap deterministic layer, a taint layer, an optional
learned layer), but it is the residual, not the authority.

\paragraph{What is the source of truth?} An append-only log. One JSON line
per event, fsynced. The model's context window, the transcript you see, the
summaries a later phase will add\,---\,all are derived from the log and can
be recomputed from it. A crash loses nothing; a replay reproduces a session
byte for byte.

\paragraph{And one more thing.} \nibb{} extends itself. The tools it uses to
build itself\,---\,run the tests, show a diff, search the tree\,---\,are written in
its own Scheme, loaded from the repository, and cross the same guard as
everything else. From the fourth week of its life, \nibb{} has been the tool
used to build \nibb{}.

\section*{Shape}

\begin{itemize}
\item \textbf{Stack.} Go host; a user-written Scheme interpreter embedded in
  it; a single static binary; a line-mode front-end now, a TUI later, and a
  core that knows nothing about either.
\item \textbf{Model.} A local Qwen served by Ollama on a workstation on the
  tailnet. No API key, no cloud dependency, and\,---\,until credentials exist\,---\,no
  secret anywhere on the build machine.
\item \textbf{Scope.} One developer, three to six months, v0. No sub-agents,
  no GUI, no embeddings, no cross-session memory. Those are decisions, not
  omissions.
\end{itemize}

\chapter{Invariants}

Five statements are load-bearing. Everything else in the design is a
consequence of them, and every phase of the plan is judged by whether it
preserves them.

\begin{enumerate}[leftmargin=2em]
\item \textbf{The log is the single source of truth.} Index and window are
  derived, revocable views; a stale summary means recompute-from-log.
\item \textbf{The judge judges intent and plans over redacted content; the
  local stack judges \emph{what}.} Jev\,---\,the external calibrated judge used
  at two seams\,---\,sees the ask, the plan, and egress \emph{situations}
  (class, destination, scope), after redaction. It never sees a secret
  value. It is never audited for the harness's own control-plane key.
  (Wording amended by \adr{19}.)
\item \textbf{Secrets are credentials, not content.} The harness holds opaque
  handles into a log-excluded encrypted keychain, uses them only at a bound
  destination, and never sees the value.
\item \textbf{The Scheme interpreter is the only Scheme$\to$Go surface; the
  guard sits on it.}
\item \textbf{Every I/O edge is guarded\,---\,by the mechanism that fits its
  substrate.} Scheme values get provenance tags; text gets redaction; egress
  gets the binding check.
\end{enumerate}

Invariant~4 is the one you will meet first in the code: \code{applyPrim} in
\file{internal/scheme/prim.go} is the only caller of any primitive's handler,
and a test records every crossing to prove it. Invariant~1 is the one that
shapes everything you will not see: there is no ``conversation state'' object
anywhere; there are events, and functions that fold them.

\chapter{Secrets are credentials}

\section{Binding, not recognition}

A secret is never suspicious in isolation. The bytes \code{sk-ant-...} sent to
\code{api.anthropic.com} are correct; the same bytes sent to
\code{collect.example.com} are exfiltration. Rightness is a \emph{binding}
question:

\begin{quote}
Pass a secret-bearing egress \emph{iff} its class is authorised for
(destination, positions, methods).
\end{quote}

The binding table is user-owned: rows of \code{(class $\to$ destination,
positions, methods)}. A bound egress passes silently and is logged. An
unbound or unknown one is denied by default (a confirm is opt-in policy)\,---\,it
never fails open. The harness does not \emph{classify} a secret; the user
\emph{binds} one, and the class is a name the user chose when registering it.

\section{Two planes}

The harness's own model key\,---\,when it has one\,---\,is \emph{control plane}:
its identity, its standing; not something it can grant or revoke, and not
something the judge audits. Every \emph{other} credential the user registers
is \emph{data plane}: bound like any other, injected only at its destination.
Under this split, other LLM providers become tools among tools: an OpenAI or
Anthropic call is an egress with a bound credential, judged exactly like an
HTTP request to anywhere else. In the bootstrap, the control plane is a local
model that needs no key at all, so the first real secret on the machine will
be the judge's own API key, in Phase~5.

\section{The content-detection ladder}

Binding is the authority, but three local, zero-retention layers watch
content, in a fixed order, each cheaper than the next and each able to stop
the escalation:

\begin{center}
\begin{tabular}{@{}lp{2.2cm}p{7.2cm}p{2.6cm}@{}}
\toprule
Layer & Acts on & Fires when & Cost \\
\midrule
L1 & text I/O & a known secret \emph{shape}: regex, entropy, prefix. Fails closed. & sub-millisecond \\
L2 & Scheme values & a value tagged at ingress as a known class heads for egress. O(1) flag at the FFI. & nothing \\
L3 & text I/O & only the class-unknown residual L1/L2 did not decide. A calibrated $P(\text{secret})$ that feeds the decision with binding, never alone. & learned; optional \\
\bottomrule
\end{tabular}
\end{center}

The deterministic layers early-terminate the learned one: if L1 or L2 already
\emph{knows}, L3 is not paid for. L2 is scoped honestly\,---\,it stops the
\emph{model's} indirect leak, not a user who can strip a tag on purpose, and
the spec rates it accordingly (0.54).

\section{The four edges}

Every path between the outside world and the two substrates\,---\,Scheme values
and the model's text context\,---\,is an edge with its own mechanism:

\begin{center}
\begin{tabular}{@{}lllp{6cm}@{}}
\toprule
Edge & From $\to$ to & Substrate & Mechanism \\
\midrule
E1 & external $\to$ Scheme value & scheme & L2 taint (provenance tag) \\
E2 & external $\to$ model context & text & L1+L3 redaction to \code{[REDACTED class]} \\
E3 & model/Scheme $\to$ external & both & binding check + L1/L3 (+ L2 flag) \\
E4 & user/OS $\to$ model context & text & L1+L3 redaction (catches pasted secrets) \\
\bottomrule
\end{tabular}
\end{center}

Redaction is a \emph{substrate controller}: it runs at every ingress to the
context, not just on API responses, so a leaked value can never be
re-emitted, even to a bound destination. The category error the design had
to resolve was thinking a taint on a value could protect the context window;
it cannot (a window is text, not a value), so the text layers run there.

\section{The judge: where it is allowed}

Jev is an external calibrated judge with no retention guarantee. The
original spec drew the line at ``never content''; \adr{19} redrew it where
it belongs, at ``never a secret value''. The question at Seam A\,---\,\emph{is
this plan good enough for the user's ask?}\,---\,has no answer without the ask
and the plan, and both are content.

So Jev is treated as what it is: an \emph{egress edge}. Its credential is a
bound data-plane credential (E3), and its payload passes the same
redaction controller as the model context (E2/E4) before it leaves. No
bespoke sanitiser; two existing mechanisms composed. What may go beyond
the ask and the plan is a user policy\,---\,\code{jev.content}: \code{none},
\code{plan} (default), \code{plan+summaries}, or \code{full}, the last
opt-in per repository because it is the level at which proprietary code is
retained elsewhere\,---\,and every call is logged with its level, size, payload
hash and the redactions it triggered.

Two seams:

\begin{description}[leftmargin=1.6em]
\item[Seam A] score a draft \emph{plan} against the \emph{ask} for quality,
  coherence and fulfilment; the model revises on a low score, at most three
  rounds.
\item[Seam C] score an ambiguous egress \emph{situation}\,---\,class,
  destination, scope\,---\,never the value.
\end{description}

Never: a secret value (enforced by redaction, not by withholding content),
control flow, its own loop, or the harness's control-plane credential. The
principle, restated: the judge judges intent and plans over \emph{redacted}
content; the local stack judges provenance and secret shape, and is the
only thing that ever sees a value.

\chapter{The loop and the guard}

The agent loop is the plain one: ask the model; run the tool it asks for,
through the guard; feed the result back; ask again. Stop on a plain answer or
a hard stop. No sub-agents in the core loop.

The loop is itself a Scheme procedure, \code{agent-turn}, so how a turn goes
can change without changing Go: plan first, check with the judge, stop a
batch of calls at the first failure. What a turn must not skip stays in Go,
behind a few primitives: windows built from the log, only the model's own
calls run, exactly one result per call, and budgets and a progress guard
that the procedure cannot catch (\adr{26}). The same procedure can route
work: narrow the context to a task and a plan, and hand it to another model
the user configured (\adr{27}).

Before each message, a plan gate: a yes/no from the judge on the abstract
intent (``is this multi-step?''), logged per message, plus a code-enforced
hard trigger for irreversible or high-stakes actions (force pushes,
recursive deletes, writes outside the tree, any unbound egress). The gate is
chosen over a pure pull model because it is auditable per message and
improves as the judge improves.

The guard is a pipeline of \emph{checkers}. Each looks at the primitive's
declared metadata and its actual arguments and returns a verdict; the
strictest wins (deny over confirm over allow). A confirm is put to the user
through a \emph{confirmer} the core implements as an event plus a wait, so
the guard never knows what a terminal is. A \emph{standing approval}\,---\,a
command prefix or a write path the user authorised ahead of time\,---\,turns a
confirm into an allow and never turns a deny into anything. Every decision
that reaches the host, and every non-allow, is logged with who asked (the
user, the model or the harness) and why.

\chapter{Memory}

Three tiers, one truth:

\begin{itemize}
\item \textbf{Disk}: the lossless log plus a session-scoped structural index.
  Source of truth.
\item \textbf{Cache}: summaries, written \emph{as log entries}, rebuildable and
  revocable, possibly stale.
\item \textbf{RAM}: the token-budgeted window, with surgical pull on demand.
\end{itemize}

v0 has no cross-session memory and no embedding index; retrieval is text
search (grep). The bootstrap ships a deliberately naive
window\,---\,elide the oldest tool results, then the oldest turns\,---\,but even it is
rebuilt from the log on every request, and it elides \emph{in place} so the
request prefix never changes and the model server's cache keeps hitting.

\chapter{Honest limits}

The spec records its own confidence, and the low numbers are the useful
ones.

\begin{enumerate}[leftmargin=2em]
\item L2 is scoped, not a security boundary. It stops the model's
  unintended leak; a user can strip a tag. The real lines are binding
  (E3), fail-closed defaults, and ingress redaction (E2/E4).
\item A bare ``confirm'' can be uninformed consent. If the class is also
  unknown, the prompt is a rubber stamp. Mitigation: always name class and
  destination; prefer requiring a binding to a bare yes.
\item L3's value is its training data: strong for known shapes, weak for a
  novel internal token\,---\,exactly the residue binding catches.
\item ``The egress surface is closed'' stands at 0.30 until redaction is a
  substrate controller and E4 is covered. Once all four edges are guarded,
  the honest statement is: no path exists for a secret \emph{value} to egress
  unless bound or explicitly re-confirmed\,---\,but ``closed'' is the user's
  affirmation, not something the harness can prove of itself.
\end{enumerate}

And in the bootstrap, before any of P3--P7 exist, there is a fifth: the
\code{shell-exec} hole. A \code{y} to \code{cat \textasciitilde/.ssh/id\_rsa}
is a \code{y}. The prompt names the path; the operator reads it.
